\documentclass[11pt,a4paper]{article}
\usepackage[utf8]{inputenc}
\usepackage[T1]{fontenc}
\usepackage{lmodern}
\usepackage[french]{babel}
\usepackage{amsmath,amssymb,mathtools}
\usepackage{icomma}
\usepackage{xcolor}
\usepackage[most]{tcolorbox}
\usepackage{enumitem}
\usepackage[a4paper,margin=2.2cm,headheight=26pt,headsep=10pt]{geometry}
\usepackage{fancyhdr}
\usepackage[hidelinks]{hyperref}
\definecolor{primary}{RGB}{222,49,99}
\definecolor{primarydark}{RGB}{150,22,55}
\definecolor{excol}{RGB}{0,135,90}
\tcbset{boxbase/.style={enhanced,breakable,boxrule=0.9pt,arc=2.5pt,
  left=3mm,right=3mm,top=2.3mm,bottom=2.3mm,fonttitle=\bfseries,coltitle=white}}
\newtcolorbox[auto counter]{corrige}[1][]{boxbase,colback=excol!4,colframe=excol,
  title={\textsc{Corrigé}~\thetcbcounter\ifx\relax#1\relax\relax\else\textmd{ --- \itshape #1}\fi}}
\newcommand{\R}{\mathbb{R}}
\newcommand{\ds}{\displaystyle}
\pagestyle{fancy}\fancyhf{}
\fancyhead[L]{\small\textcolor{primarydark}{\textbf{Thème 2} — Logarithmes}}
\fancyhead[R]{\small\textcolor{primarydark}{\textsc{Corrigé — Facile}}}
\fancyfoot[C]{\small\thepage}
\renewcommand{\headrulewidth}{0.8pt}
\begin{document}

\begin{center}
{\LARGE\bfseries\color{primarydark} Niveau Facile — Corrigé}\\[2pt]
{\color{primary}\rule{0.45\textwidth}{1.2pt}}
\end{center}
\vspace{1mm}

\begin{corrige}[la définition]
On cherche l'exposant à donner à la base.
\[ \text{a) }\log_2 8=3\ (2^3=8);\quad \text{b) }\log_3 81=4\ (3^4=81);\quad \text{c) }\log_{10}1000=3;\quad
   \text{d) }\log_5 1=0;\quad \text{e) }\ln e=1. \]
\end{corrige}

\begin{corrige}[développer]
\[ \text{a) }\log_a x+\log_a y;\quad \text{b) }\log_a x-\log_a y;\quad \text{c) }3\log_a x;\quad
   \text{d) }2\ln x;\quad \text{e) }\log_a(x^{2}y)=2\log_a x+\log_a y. \]
\end{corrige}

\begin{corrige}[valeurs remarquables]
\[ \text{a) }1;\quad \text{b) }\ln e^{4}=4\ln e=4;\quad \text{c) }0;\quad
   \text{d) }\ln\tfrac1e=-\ln e=-1;\quad \text{e) }\log_{10}0{,}01=\log_{10}10^{-2}=-2. \]
\end{corrige}

\begin{corrige}[isoler l'exposant]
\begin{align*}
&\text{a) }x=\ln 7; &&\text{b) }x=\log_{10}100=2; &&\text{c) }2x=5\Rightarrow x=\tfrac52;\\
&\text{d) }-x=\ln 3\Rightarrow x=-\ln 3; &&\text{e) }x=\log_{10}5=\log 5. &&
\end{align*}
\end{corrige}

\begin{corrige}[le piège des carrés]
\begin{itemize}[leftmargin=1.4em,topsep=1pt,itemsep=1pt]
  \item a) \textbf{Vrai} : $\ln(x^{2})=2\ln x$ (propriété du log d'une puissance).
  \item b) \textbf{Faux} : $\ln^{2}x=(\ln x)^{2}\neq 2\ln x$ en général.
  \item c) \textbf{Vrai} : $\ln^{2}x$ est justement une \emph{notation} pour $(\ln x)^{2}$.
  \item d) \textbf{Faux} : $\ln(x^{2})=2\ln x$, alors que $\ln^{2}x=(\ln x)^{2}$ : ce sont deux objets différents.
\end{itemize}
\end{corrige}

\end{document}
